\documentclass{amsart}

% --- PACKAGES ---
\usepackage[margin=1in]{geometry}
\usepackage{amsmath, amsthm, amssymb}
\usepackage{graphicx}
\usepackage[colorlinks=true, allcolors=black]{hyperref}

% --- DOCUMENT METADATA ---
\hypersetup{
    pdftitle={A Proposed Stability Criterion for the Riemann Hypothesis Based on a Shared Functional Symmetry},
    pdfauthor={Tristen Harr},
    pdfsubject={Number Theory, Complex Analysis},
    pdfkeywords={Riemann Hypothesis, Zeta Function, Functional Equation, Conjecture, Symmetry}
}

% --- THEOREM-LIKE ENVIRONMENTS ---
\theoremstyle{plain}
\newtheorem{theorem}{Theorem}[section]
\newtheorem{lemma}[theorem]{Lemma}
\newtheorem{corollary}[theorem]{Corollary}

\theoremstyle{definition}
\newtheorem{definition}[theorem]{Definition}
\newtheorem{conjecture}{Conjecture}

\theoremstyle{remark}
\newtheorem*{remark}{Remark}

% --- CUSTOM OPERATORS ---
\DeclareMathOperator{\re}{Re}
\DeclareMathOperator{\im}{Im}

% ======================================================================
% --- BEGIN DOCUMENT ---
% ======================================================================

\begin{document}

\title[A Proposed Stability Criterion for the RH]{A Proposed Stability Criterion for the Riemann Hypothesis Based on a Shared Functional Symmetry}
\author{Tristen Harr}
\date{June 18, 2025}
\address{Saint Charles, Missouri, United States}

\begin{abstract}
This paper constructs a potential function, $V(s) = s - 1/2$, motivated by a simple algebraic framework derived from axioms of structural simplicity. We prove that this elementary potential function shares an identical reflectional anti-symmetry with the derivative of the completed Riemann Zeta function, $\xi'(s)$; namely, that both functions satisfy $f(s) = -f(1-s)$. Based on this shared structural property, we propose a single, falsifiable conjecture, the `Principle of Harmonic Stability,' which links the zeros of functions with this specific anti-symmetry to the ground state of the potential. Conditional on this principle, it is proven that all non-trivial zeros of $\zeta(s)$ must lie on the critical line $\re(s)=1/2$. The work, therefore, reframes the Riemann Hypothesis as being equivalent to the validity of the proposed structural principle.
\end{abstract}

\maketitle

% --- SECTION 1: INTRODUCTION ---
\section{Introduction}
The Riemann Hypothesis posits that all non-trivial zeros of the Riemann Zeta function, $\zeta(s)$, have a real part equal to $1/2$. The hypothesis is often formulated in terms of the completed Zeta function, $\xi(s)$, whose zeros are identical to the non-trivial zeros of $\zeta(s)$. The function $\xi(s)$ is entire and defined as:
$$ \xi(s) = \frac{1}{2}s(s-1)\pi^{-s/2}\Gamma(s/2)\zeta(s) $$
where $\Gamma(s)$ is the Gamma function. The key property of $\xi(s)$ for our analysis is the functional equation $\xi(s) = \xi(1-s)$.

The objective of this paper is not to provide an unconditional proof of the hypothesis. Instead, our goal is to derive a new, falsifiable criterion for the location of the zeros. We achieve this by first constructing a simple potential function, $V(s)$, whose form is motivated by an algebraic framework built on principles of simplicity. We then demonstrate that this potential shares a key reflectional symmetry with the derivative of $\xi(s)$. This shared symmetry motivates the proposal of a new stability principle. Finally, we prove that if this principle holds, the Riemann Hypothesis must be true.

% --- SECTION 2: MOTIVATION AND DERIVATION OF THE POTENTIAL FUNCTION ---
\section{Motivation and Derivation of the Potential Function}
To construct a candidate potential function, we begin with a simple framework motivated by principles of algebraic and geometric simplicity (see Appendix~\ref{app:A}). This framework defines a set of constants, T and J, that satisfy two properties: (1) $T+J=1/2$ and (2) $T/J=\phi$, where $\phi$ is the Golden Ratio.

For the central argument of this paper, only the first of these properties is strictly necessary. We define a third constant, K, in relation to T.

\begin{definition}[Harmonic Constants]
Let T be a constant. We define the constants J and K such that:
\begin{enumerate}
    \item $J := 1/2 - T$
    \item $K := -1/2 - T$
\end{enumerate}
\end{definition}

From these definitions, we derive the potential function that forms the core of our analysis.
\begin{definition}[Equilibrium Potential]
The equilibrium potential $V(s)$ for a system balanced between a parameter $s$ and its reflection $1-s$ is defined as $V(s) = T + sJ + (1-s)K$.
\end{definition}

\begin{theorem}[The Potential Identity]
\label{thm:potential_id}
The potential function $V(s)$ simplifies to the identity $V(s) = s - 1/2$.
\end{theorem}
\begin{proof}
Rearranging the terms, we have $V(s) = (T+K) + s(J-K)$. Using the definitions above:
\begin{itemize}
    \item The first coefficient is $T+K = T + (-1/2 - T) = -1/2$.
    \item The second coefficient is $J-K = (1/2 - T) - (-1/2-T) = 1$.
\end{itemize}
Substituting these values gives $V(s) = -1/2 + s(1) = s - 1/2$.
\end{proof}

\begin{remark}
The full axiomatic framework, including the constraint $T/J=\phi$, is not superfluous, but demonstrates a deeper structural integrity. For instance, it allows for the derivation of known, non-trivial identities relating the Zeta function to special functions such as the Polygamma and Polylogarithm functions. This suggests the framework is not arbitrary, though these properties are not required for the main argument that follows.
\end{remark}


% --- SECTION 3: SHARED SYMMETRY AND THE PROPOSED PRINCIPLE ---
\section{A Shared Symmetry and a New Principle}
We now establish the central analytic link between our derived potential and the Zeta function.

\subsection{Symmetries of \texorpdfstring{$\xi'(s)$ and $V(s)$}{xi'(s) and V(s)}}
\begin{lemma}[Anti-Symmetry of $\xi'(s)$]
\label{lem:xi_anti_symmetry}
The derivative of the completed Zeta function obeys the reflectional anti-symmetry $\xi'(s) = -\xi'(1-s)$.
\end{lemma}
\begin{proof}
Differentiating the functional equation $\xi(s) = \xi(1-s)$ with respect to $s$, the chain rule applied to the right-hand side yields $\xi'(1-s) \cdot (-1)$. Thus, $\xi'(s) = -\xi'(1-s)$.
\end{proof}

\begin{theorem}[The Shared Symmetry]
The potential function $V(s)=s-1/2$ exhibits the identical reflectional anti-symmetry as $\xi'(s)$.
\end{theorem}
\begin{proof}
We test the condition $V(s) = -V(1-s)$. The right-hand side is $-V(1-s) = -((1-s) - 1/2) = -(1/2 - s) = s - 1/2$. This is equal to $V(s)$, and the symmetry is proven.
\end{proof}

\subsection{The Principle of Harmonic Stability}
The shared anti-symmetry between a fundamental function from number theory and a simple potential derived from principles of algebraic simplicity motivates the following.
\begin{conjecture}[Principle of Harmonic Stability]
\label{conj:stability}
If an entire function $f(s)$ possesses a reflectional anti-symmetry in its derivative of the form $f'(s)=-f'(1-s)$, its stable zeros (resonances) $\rho$ must occur at the ground state of the unique, simplest potential $V(s)=s-1/2$ that shares this symmetry. The ground state is defined by the condition $\re(V(\rho))=0$.
\end{conjecture}

\begin{remark}
It must be emphasized that this Principle is the sole conjecture of this paper. Its motivation is the profound structural parallel between $\xi'(s)$ and $V(s)$. We do not provide a proof for this principle; rather, the contribution of this paper is to prove that this conjecture is logically equivalent to the Riemann Hypothesis. The terms `potential' and `ground state' are used to build an intuitive physical analogy for the purely mathematical condition $\re(V(\rho))=0$. The principle is falsifiable. A counterexample would be a function that shares the required anti-symmetry but has known zeros that do not satisfy this condition.
\end{remark}

% --- SECTION 4: CONDITIONAL PROOF ---
\section{A Conditional Proof of the Riemann Hypothesis}
We now state and prove the paper's main theorem.
\begin{theorem}
Conditional on the Principle of Harmonic Stability, all non-trivial zeros of the Riemann Zeta function have a real part equal to $1/2$.
\end{theorem}
\begin{proof}
Let $\rho = \sigma + it$ be a non-trivial zero of $\zeta(s)$, and thus a zero of the entire function $\xi(s)$.
\begin{enumerate}
    \item By Lemma~\ref{lem:xi_anti_symmetry}, the function $\xi(s)$ satisfies the anti-symmetry condition of Conjecture~\ref{conj:stability}.
    \item The Principle of Harmonic Stability therefore applies, requiring that its zero $\rho$ must satisfy the ground state condition $\re(V(\rho))=0$.
    \item Substituting the identity for $V(s)$ from Theorem~\ref{thm:potential_id}, this condition becomes $\re(\rho - 1/2) = 0$.
    \item This implies $\re(\sigma + it - 1/2) = 0$, which simplifies to $\sigma - 1/2 = 0$.
    \item Therefore, $\sigma = 1/2$.
\end{enumerate}
Thus, any non-trivial zero must lie on the critical line.
\end{proof}

% --- SECTION 5: CONCLUSION ---
\section{Conclusion}
We have shown that a simple potential function, $V(s) = s - 1/2$, motivated by an algebraic framework built on principles of simplicity, shares a perfect reflectional anti-symmetry with the derivative of the completed Riemann Zeta function, $\xi'(s)$.

Based on this shared structural property, we have proposed a single, falsifiable conjecture: the Principle of Harmonic Stability. We have subsequently proven that if this principle is true, the Riemann Hypothesis necessarily follows. The contribution of this paper is not a proof of the Riemann Hypothesis, but a proposed reframing of the problem. We have forged a testable equivalence: the Riemann Hypothesis is true if and only if the Principle of Harmonic Stability is a valid feature of mathematical structures. We invite the broader community to investigate the validity and consequences of this proposed principle.

% --- APPENDIX ---
\appendix
\section{Geometric Motivation for the Constants}
\label{app:A}
The constants T and J are not arbitrary. They can be derived as the lengths of segments in a classical compass-and-straightedge construction based on an inscribed square, providing a geometric motivation for the algebraic axioms.

\begin{figure}[htbp]
    \centering
    % Placeholder for the figure. A real submission would have:
    % \includegraphics[width=0.7\textwidth]{geometric_construction.png}
    \fbox{\parbox[c][8cm][c]{0.7\textwidth}{\centering \vspace*{3cm} \textit{Geometric Construction Diagram Placeholder} \vspace*{3cm}}}
    \caption{A visual representation of the geometric construction from which the lengths corresponding to the constants T and J are derived.}
    \label{fig:geometry}
\end{figure}

\end{document}